\documentclass[11pt,a4paper]{article}
\usepackage[utf8]{inputenc}
\usepackage[T1]{fontenc}
\usepackage[margin=1in]{geometry}
\usepackage{amsmath,amssymb,amsfonts,amsthm}
\usepackage{booktabs}
\usepackage{microtype}
\usepackage{titlesec}
\usepackage{hyperref}
\usepackage{cleveref}
\usepackage{siunitx}
\usepackage{enumitem}
\usepackage{float}
\usepackage{graphicx}
\usepackage{longtable}

\titleformat{\section}{\large\bfseries}{\thesection}{1em}{}
\titleformat{\subsection}{\normalsize\bfseries}{\thesubsection}{1em}{}

\hypersetup{
    colorlinks=true,
    linkcolor=blue,
    citecolor=blue,
    urlcolor=blue,
    pdfauthor={Michael Lehmann},
    pdftitle={RDCC – Speculative Extensions \& Conceptual Companions II}
}
    
\sisetup{range-phrase=--, range-units=single, per-mode=symbol}

\title{\textbf{RDCC – Speculative Extensions \& Conceptual Companions II}}
\author{Michael Lehmann \\
        Independent Researcher \\
        \texttt{mi.lehmann@gmx.de}}
\date{March 2026}

\begin{document}
\maketitle
\vspace{1em}
\begin{center}

\textbf{This document is not part of the predictive RDCC framework but a conceptual interpretation!}

\vspace{1em}

\textbf{These analogies do not validate RDCC, but demonstrate that its underlying structure is deeply rooted in known physics.}
\end{center}
\vspace{1em}

\begin{abstract}
This document presents a collection of structural, interpretative, and speculative extensions of the Relaxation-Driven Cyclic Cosmology (RDCC) framework. These ideas do not modify the core dynamics of RDCC and introduce no new fields, interactions, or predictive mechanisms. Instead, they explore how global CPT symmetry, sectoral projection, entanglement structure, and the limits of description shape our understanding of time, information, and high-curvature regimes. The aim is to provide a conceptual companion to RDCC, clarifying how its structural principles naturally extend into broader interpretative domains without compromising its minimalist philosophy.
\end{abstract}

\section{Introduction}

Relaxation-Driven Cyclic Cosmology (RDCC) is built on a single structural postulate: the global quantum state of the Universe is a CPT eigenstate. This principle leads to a bipartite Hilbert-space factorization, sectoral arrows of time, and a nonsingular bounce within standard Einstein gravity. The framework is intentionally minimal: no new fields, no modified gravity, and only one dimensionless parameter, the relaxation strength $\alpha$.

The present document collects conceptual extensions that arise naturally from this structural foundation. These ideas are not part of the RDCC core model and do not aim to provide new predictions. Instead, they illuminate how RDCC's dualistic structure offers a coherent interpretative lens for entanglement, black hole interiors, causal limits, and high-curvature regimes.

All considerations remain strictly non-dynamical and compatible with the minimalist spirit of RDCC.

\section{Global CPT Symmetry and Sectoral Projection}

\subsection{Global Non-Separability}

The global CPT-symmetric state is non-separable. Sectoral descriptions arise only after projection onto one branch of the bipartite Hilbert space. This projection induces three key features of sectoral physics:

1. sectoral time orientation  
2. sectoral entropy gradients  
3. sectoral causal structure  

Globally, none of these exist. The Universe is timeless and symmetric at the fixpoint, and sectoral distinctions emerge only through projection.

\subsection{Entanglement as Structural Relation}

Entanglement is interpreted as a structural signature of global non-separability. Sectoral observers perceive correlations as nonlocal, but globally they are simply relational constraints within a CPT-symmetric state.

This perspective aligns with thermofield-double constructions and ER=EPR-type interpretations, without requiring geometric identifications or additional dynamical assumptions. Entanglement thus appears as a boundary effect of sectoral description rather than a fundamental nonlocal mechanism.

\subsection{Purification Principle}

Any mixed state $\rho_+$ on a single sector can always be purified to a pure state on the bipartite Hilbert space $\mathcal{H}_+ \otimes \mathcal{H}_-$. By the purification theorem there exists a pure state $|\Psi\rangle$ such that
\[
\rho_+ = \operatorname{Tr}_- \big( |\Psi\rangle\langle\Psi| \big).
\]
The global state $|\Psi\rangle$ is pure, timeless, and CPT-symmetric, while the sectoral reduced state appears mixed and time-directed. This is the most general mathematical realization of the RDCC structure: every sectoral “reality” is the projection of a globally non-separable pure state.

\subsection{Thermofield Double States}

The thermofield double (TFD) state provides an almost identical realization of the global CPT eigenstate of RDCC:
\[
|\text{TFD}\rangle = \frac{1}{\sqrt{Z(\beta)}} \sum_n e^{-\beta E_n/2} |n\rangle_+ \otimes |n\rangle_-.
\]
The global TFD state is pure and symmetric, yet each individual sector appears to its observer as a thermal state with temperature $1/\beta$. This is precisely the RDCC picture: a single global CPT eigenstate whose sectoral projections carry thermodynamic arrows and entropy gradients.

\section{Fixpoints and the Limits of Sectoral Description}

\subsection{The Global CPT Fixpoint}

The cosmological bounce corresponds to a global fixpoint where:

1. entropy gradients vanish  
2. time orientation disappears  
3. the relaxation parameter $\alpha$ approaches zero and suppresses sectoral distinction  

This fixpoint is not a dynamical event but a structural boundary of description. It marks the regime in which sectoral notions of time, causality, and locality lose operational meaning.

\subsection{Black Holes as Local Fixpoints}

Black hole interiors can be interpreted as local analogues of the global fixpoint. At the horizon, the effective Hubble parameter vanishes, mirroring the bounce condition. Information is preserved structurally through global non-separability rather than dynamically recovered.

This reframes the information paradox as a limit of sectoral description rather than a failure of unitarity. Black holes thus serve as local windows into the same structural symmetry that governs the cosmological bounce.

\subsection{Spacetime Bipartition: The Unruh Effect}

The Minkowski vacuum $|0_M\rangle$ is a pure, Poincaré-invariant global state. It factorizes into a bipartite entangled state between the left and right Rindler wedges:
\[
|0_M\rangle \sim \sum_i c_i |i\rangle_L \otimes |i\rangle_R.
\]
An accelerated observer in one wedge experiences a thermal bath (Unruh temperature $T = a/2\pi$). The global state is timeless and symmetric, while each sectoral projection produces a thermodynamic arrow and entropy. This is a direct geometric realization of the RDCC mirror-sector structure.

\subsection{Hartle--Hawking State for Eternal Black Holes}

The Hartle--Hawking–Israel state on the Kruskal manifold is a pure global state that is maximally entangled across the event horizon. The two asymptotic regions (exterior and “mirror” interior) form a bipartite CPT-symmetric system. At the horizon (the local fixpoint) sectoral distinctions vanish, exactly as at the cosmological bounce. This provides the precise quantum-field-theoretic realization of the local fixpoint analogy already discussed.

\subsection{Critical Points and Spontaneous Symmetry Breaking (Analogy)}

A useful structural analogy for the CPT fixpoint in RDCC is provided by critical phenomena in condensed-matter systems. Consider a ferromagnet near its critical temperature $T_{c}$. Above $T_{c}$ the system is symmetric and exhibits no preferred magnetization direction. Below $T_{c}$ the effective Landau potential
\begin{equation}
    V(\phi) = a\,\phi^{2} + b\,\phi^{4}
\end{equation}
develops two degenerate minima at $\phi>0$ and $\phi<0$. The system must choose one branch, and this choice defines the macroscopic order parameter.

This structure mirrors the RDCC picture. At the global CPT fixpoint the Universe is symmetric, timeless, and without an entropy gradient. After sectoral projection, the two CPT-related branches correspond to the two degenerate minima: one with a forward-pointing thermodynamic arrow, the other with a reversed arrow. The fixpoint itself is not a dynamical transition but a structural boundary at which sectoral distinctions lose operational meaning.

The analogy highlights that the emergence of a time orientation in each sector is not a fundamental asymmetry but a spontaneous selection of one branch of a globally symmetric state, much like the spontaneous magnetization of a ferromagnet below $T_{c}$.

\section{Structural Analogies: Condensed Matter and RDCC}

Several phenomena in condensed-matter physics exhibit structural features that closely mirror the dualistic architecture of RDCC. These analogies do not imply dynamical equivalence; rather, they highlight how global symmetry, local projection, and the emergence of sectoral structure arise in diverse physical systems.

\subsection{BCS Superconductivity: Global Coherence vs. Sectoral Description}

In BCS theory, electrons form Cooper pairs that condense into a single macroscopic quantum state. The global condensate is coherent and symmetric, while the sectoral description in terms of individual electrons appears asymmetric and dynamical. The transition at the critical temperature marks a boundary where the sectoral description loses validity and global coherence dominates.

This mirrors the RDCC distinction between the globally CPT-symmetric state and the sectoral projection into visible and mirror branches. In both cases, the global description is fundamental, while the sectoral picture is emergent and limited.

\subsection{Kibble--Zurek Mechanism: Fixpoints and Branch Selection}

When a system is driven through a critical point, it cannot choose a symmetry-breaking branch uniformly. Instead, domains form, each selecting one of the degenerate minima. The critical point itself is a structural boundary where the correlation length diverges and the usual sectoral description breaks down.

RDCC exhibits an analogous structure: the global CPT fixpoint is symmetric and timeless, while sectoral time orientations emerge only after projection. The fixpoint is not a dynamical transition but a boundary of description, much like the critical point in the Kibble--Zurek mechanism.

\subsection{Ising Duality: Local vs. Global Descriptions}

The Ising model admits two complementary descriptions: a local spin picture and a dual description in terms of domain walls. These two views are mathematically equivalent but highlight different structural aspects of the same system. Local asymmetries coexist with global symmetry.

RDCC exhibits a similar duality. The sectoral description features time orientation, entropy gradients, and causal structure, while the global CPT-symmetric state is timeless and non-separable. The apparent asymmetries of sectoral physics arise from projection, not from fundamental breaking of symmetry.

\subsection{Holographic Principle and AdS/CFT Correspondence}

The holographic principle, most explicitly realized in the AdS/CFT correspondence, provides a striking structural parallel to RDCC. In the global, higher-dimensional bulk description the quantum state is non-separable and contains all information in a single, CPT-symmetric wave function. Sectoral observers, however, access only the lower-dimensional boundary theory (the “visible sector”), from which local spacetime geometry, causality, and entropy gradients emerge through projection. The horizon or boundary acts as a structural fixpoint where bulk and boundary descriptions become equivalent and sectoral notions of locality lose operational meaning. This mirrors the RDCC picture in which the global CPT eigenstate is projected onto visible and mirror sectors, with the cosmological bounce serving as the fixpoint at which sectoral distinctions vanish.

\subsection{Wheeler--DeWitt Equation and Timeless Quantum Cosmology}

In canonical quantum gravity the Wheeler--DeWitt equation describes a global, timeless wave functional $\Psi[g, \phi]$ of the entire universe. Time, entropy gradients, and causal structure appear only in a semiclassical, sectoral approximation. The global state remains CPT-symmetric and non-separable; the emergence of a thermodynamic arrow is a projection effect exactly analogous to the RDCC sectoral arrows that arise after the bounce. The cosmological fixpoint ($H = 0$, $\alpha \to 0$) corresponds to the regime where the Wheeler--DeWitt description becomes exact and the usual sectoral time coordinate disappears, reinforcing that time is not fundamental but a boundary phenomenon of sectoral description.

\subsection{Decoherence and Environmental Projection}

Decoherence arises when a globally entangled quantum state is traced over an “environment” (the mirror sector in RDCC language). The reduced density matrix of the visible sector then appears classical and time-directed, even though the global state remains pure and unitary. The pointer basis and apparent collapse are projection effects, not fundamental non-unitary processes. This structure is identical to the RDCC mechanism: the global CPT eigenstate is non-separable; only after sectoral projection does a thermodynamic arrow and classical spacetime geometry emerge. The bounce, where maximal entanglement is restored ($\alpha \to 0$), corresponds to the point of perfect global coherence before any sectoral decoherence sets in.

\subsection{Quantum Teleportation and Relational Constraints}

Quantum teleportation with Bell pairs illustrates the same relational structure. Globally, the entangled state is non-separable and timeless; information is encoded in the full CPT-symmetric correlations. Only after sectoral measurement and a classical channel does the transfer appear “instantaneous” to a local observer. This apparent non-locality is not a physical influence but a projection effect of the global non-separability—precisely as RDCC interprets Bell correlations and the gravitational shadow of the mirror sector. The protocol succeeds because the global state already contains the relational constraint; the sectoral description merely makes this constraint operationally accessible.

\subsection{ER=EPR as Relational Projection (Deepening)}

ER=EPR-type correlations receive a natural structural interpretation within RDCC. Wormhole-like geometries are not required; the entanglement is the direct manifestation of global non-separability projected onto a sectoral description. The fixpoint (bounce or horizon) enforces the boundary at which such correlations become maximal while sectoral separation is still maintained.

\subsection{Many-Worlds as Multi-Sectoral Projection}

The Many-Worlds interpretation can be viewed as a maximal sectoral projection: each branch of the universal wave function corresponds to one possible realization of the global CPT-symmetric state. The global state itself remains single, non-separable and timeless; the apparent multiplicity of worlds is simply the result of restricting attention to one branch at a time. This aligns exactly with the RDCC picture in which the visible and mirror sectors are two complementary projections of the same global eigenstate.

\subsection{Black-Hole Interiors as Local CPT Fixpoints (Extended Analogy)}

Black-hole interiors constitute a local realization of the global CPT fixpoint. Inside the horizon the radial coordinate effectively becomes timelike, and the interior geometry can be regarded as the CPT-conjugate continuation of the exterior region. The would-be singularity or firewall is reinterpreted as the point at which sectoral description breaks down completely—mirroring the cosmological bounce where $\alpha\to 0$ and sectoral distinctions vanish. Information is not lost but remains encoded in the global non-separable CPT eigenstate; a sectoral observer experiences only an apparent loss because they are restricted to one branch of the Hilbert space. This local–global correspondence strengthens the structural identification between black-hole horizons and the cosmological bounce without requiring new dynamical mechanisms or UV completions. The black-hole interior thus serves as a microscopic laboratory for the same projection effects that govern the entire cyclic cosmology.

\subsection{Extended Summary}

These analogies illustrate that the dualistic structure of RDCC is not exotic but reflects a broader pattern throughout physics: global symmetry combined with sectoral projection, fixpoints as boundaries of description, and emergent asymmetry arising from structurally equivalent branches. From condensed-matter systems through holography, timeless quantum cosmology, decoherence, teleportation, ER=EPR, Many-Worlds and black-hole interiors, the same mechanism appears again and again. RDCC extends this universal pattern to cosmology, interpreting the visible and mirror sectors as complementary projections of a globally symmetric CPT eigenstate.

\section{Light Speed as a Sectoral Coherence Boundary}

\subsection{c as Structural Limit}

Within the RDCC perspective, the speed of light is not interpreted as a dynamical parameter or a property of a specific field. Instead, it functions as a structural boundary that stabilizes causal and relational coherence within a time-directed sector.

Massive systems may asymptotically approach the speed of light but never reach it. This asymptotic behavior indicates that c represents a limit of sectoral description rather than a realizable physical state. Beyond this limit, sectoral notions of simultaneity, locality, and causal ordering lose operational meaning.

\subsection{Analogy with the Global Fixpoint}

The speed of light and the global CPT fixpoint share a deep structural analogy:

1. both are unreachable limits from within a sector  
2. both mark boundaries where sectoral description breaks down  
3. both stabilize the relational structure of physical laws  

The speed of light governs the coherence of local causal relations, while the global fixpoint governs the coherence of temporal and entropic relations across the full CPT-symmetric state. In this sense, c can be viewed as a sectoral projection of the global symmetry encoded at the fixpoint.

\section{Quantum Gravity as Projection in CPT-Symmetric Cycles}

\subsection{High-Curvature Regimes as Projection Effects}

Phenomena typically associated with quantum gravity, such as strong curvature effects, loop corrections, or mode mixing, can be interpreted as sectoral manifestations of approaching the global CPT fixpoint. In this view, high-curvature regimes do not require new ultraviolet degrees of freedom. Instead, they reflect the structural suppression of sectoral distinctions as the system approaches maximal symmetry. These considerations are further illuminated by the holographic principle, the Wheeler--DeWitt equation and the extended black-hole interior analogy discussed in the structural analogies of Section 4.

The global CPT structure remains intact even when classical general relativity becomes insufficient as a sectoral description. This provides a conceptual bridge between classical and quantum-gravitational regimes without introducing new dynamical mechanisms.

\subsection{Entanglement in Strong Fields}

Near horizons or near the cosmological bounce, entanglement across sectors becomes more pronounced. This offers a structural interpretation of ER=EPR-type correlations: entanglement is not a fundamental nonlocal interaction but a projection of global non-separability onto a sectoral description.

The fixpoint acts as a structural boundary that preserves correlations while enforcing sectoral separation. This perspective avoids the need for geometric wormholes or additional quantum-gravity constructs.

\section{Mirror Photons and Gravitational Shadow}

\subsection{Global and Sectoral Fields}

Photons and gravitons are global fields in the RDCC framework. However, their quanta are strictly sectoral: visible photons couple only to visible-sector charges, and mirror photons couple only to mirror-sector charges. As a result, mirror light is electromagnetically invisible to us.

Despite this, mirror photons carry energy and momentum, and all forms of energy gravitate. The shared spacetime geometry therefore encodes the combined energetic contributions of both sectors.

\subsection{Dark Matter as Gravitational Shadow}

The gravitational influence of mirror photons and mirror matter appears to visible-sector observers as dark matter. This influence is a compressed and indirect form of information about the mirror sector. We do not observe mirror photons themselves, nor their frequencies or directions, but only the curvature they induce.

This interpretation requires no new particle species beyond those already implied by the CPT-symmetric two-sector structure. Dark matter emerges naturally as the gravitational shadow of the mirror sector.

\section{Observational Windows}

Although the ideas presented in this document are conceptual and non-dynamical, they suggest several potential observational windows:

1. low-frequency gravitational-wave tails accessible to pulsar timing arrays  
2. subtle modulations in the 21-cm signal that may reflect sectoral correlations  
3. correlations between the relaxation parameter alpha and late-time structure growth  

These possibilities are not predictions of RDCC but interpretative avenues that may guide future exploration.

\section{Limitations}

The extensions discussed in this document are qualitative, structural, and intentionally non-predictive. They do not constitute a theory of quantum gravity, nor do they introduce new fields or mechanisms. Their purpose is to clarify how RDCC's structural principles naturally extend into broader conceptual domains.

Open questions include:

1. the detailed behavior of entanglement near the fixpoint  
2. the backreaction of sectoral projection on bounce dynamics  
3. the embedding of the fixpoint structure in more complete quantum frameworks  

Addressing these questions lies beyond the scope of the present companion.

\section{RDCC and Microphysical Puzzles}

\subsection{Overview}

The structural principles of RDCC provide a perspective that does not modify quantum mechanics or general relativity, but reframes several long-standing microphysical puzzles. These puzzles arise when sectoral descriptions are applied to a globally CPT-symmetric, non-separable state. RDCC does not claim to solve these problems dynamically. Instead, it clarifies how many apparent paradoxes originate from the limits of sectoral projection.

\subsection{The Problem of Time}

In quantum mechanics, time is an external parameter. In general relativity, time is dynamical. In canonical quantum gravity, time disappears entirely. RDCC offers a structural reinterpretation: time is a sectoral construct that emerges only after projection onto one branch of the global CPT-symmetric state. Globally, the Universe is timeless. The problem of time is therefore not a physical inconsistency but a mismatch between global and sectoral descriptions.

\subsection{The Thermodynamic Arrow of Time}

The microscopic laws of physics are time-reversal symmetric, yet macroscopic systems exhibit a thermodynamic arrow. In RDCC, this arrow is not fundamental. It arises because each sector inherits a time orientation from the projection of the global state. The relation
\[
S_{+}(t) = S_{-}(-t)
\]
expresses that entropy gradients are sectoral, not global. The arrow of time is thus a perspectival effect rather than a fundamental asymmetry.

\subsection{The Measurement Problem}

Quantum measurement appears to require a non-unitary collapse, in tension with global unitarity. RDCC reframes this tension structurally. The global state remains pure and unitary, while sectoral observers access only reduced states obtained by projection. Apparent collapse corresponds to the loss of global information under sectoral restriction. This interpretation is compatible with relational and information-theoretic approaches to quantum mechanics.

\subsection{Bell Nonlocality}

Bell correlations appear nonlocal within a sector, but RDCC treats them as relational constraints within a globally non-separable state. Nonlocality is therefore not a physical influence but a limitation of sectoral description. The global CPT-symmetric state is local and timeless; nonlocality emerges only after projection.

\subsection{The Role of the Speed of Light}

The constancy and universality of the speed of light are structural features of sectoral coherence. RDCC interprets c as a boundary of sectoral description rather than a dynamical parameter. Just as the global fixpoint marks the limit of temporal description, c marks the limit of causal description. Both limits arise from the same structural symmetry.

\subsection{The Weakness of Gravity}

The hierarchy between gravitational and gauge interactions remains unexplained in standard frameworks. RDCC does not introduce new mechanisms but offers a structural perspective: gravity is a global field shared by both sectors, while gauge interactions are sectoral. The relative weakness of gravity reflects its inter-sectoral character rather than a dynamical hierarchy.

\subsection{Dark Matter as Gravitational Shadow}

RDCC interprets dark matter as the gravitational imprint of the mirror sector. Mirror photons and mirror matter gravitate but do not interact electromagnetically with the visible sector. This provides a minimal and structurally natural explanation for dark matter without introducing new particle species beyond those implied by CPT symmetry.

\subsection{The Information Paradox}

Black hole evaporation appears to violate unitarity. RDCC reframes this paradox by interpreting black hole interiors as local fixpoints analogous to the cosmological bounce. Information is preserved globally through non-separability, while sectoral observers experience an apparent loss due to projection. The paradox arises from the limits of sectoral description, not from a failure of fundamental physics.

\subsection{Vacuum Energy and the Cosmological Constant}

The smallness of the observed cosmological constant remains a major open problem. RDCC does not provide a dynamical solution but offers a structural reinterpretation: vacuum energy is a sectoral quantity, while the global state is CPT-symmetric and timeless. The observed value may therefore reflect a projection effect rather than a fundamental parameter of the global state.

\subsection{Linearity of Quantum Mechanics}

The linearity of quantum mechanics is a foundational assumption with no deeper explanation. RDCC suggests that linearity may arise from the projection of a globally non-separable state onto sectoral degrees of freedom. This perspective aligns with relational and information-based interpretations of quantum theory.

\subsection{Decoherence, Higgs Symmetry Breaking, Quantum Teleportation, Many-Worlds and Black-Hole Interiors}

Decoherence, the Higgs mechanism, quantum teleportation, the Many-Worlds interpretation and the extended black-hole interior analogy (already discussed in the structural analogies of Section 4) are further examples of projection effects. They reinforce that many apparent quantum-mechanical peculiarities—classical emergence, symmetry breaking, relational information transfer, the apparent multiplicity of worlds and the information paradox—are natural consequences of restricting a globally non-separable CPT-symmetric state to a single sector.

\section{Gravitation as the Global Bridge}
\textbf{The Only Interaction That Sees Both Sectors}

In the RDCC framework gravity occupies a unique structural position: it is the \emph{sole interaction that couples to the global CPT-symmetric state rather than to a single sector}. All other forces (electromagnetic, weak, strong) are strictly sectoral---they act only within $\mathcal{H}_+$ or $\mathcal{H}_-$. Gravity, by contrast, couples to the \emph{sum} of the two energy-momentum tensors:
\begin{equation}
\mathcal{L}_{\rm int} = \frac{1}{M_*^2} T^{(+)}_{\mu\nu} T^{(-)\mu\nu},
\label{eq:lint}
\end{equation}
exactly as introduced in the core RDCC paper. This operator is the microscopic origin of the relaxation parameter $\alpha$ and simultaneously the reason why gravity appears so weak to sectoral observers.

\subsection{The Effective Weakness of Gravity}
A visible-sector observer measures an effective Newton constant
\[
G_{\rm eff} = \frac{G_N}{1 + \rho_- / \rho_+},
\]
where $\rho_-/\rho_+ \approx x^4 \approx 0.06$ for the dynamically generated temperature ratio $x \approx 0.5$. The mirror sector contributes an invisible ``gravitational shadow'' that dilutes the apparent strength of gravity without introducing new particles or scales. Gravity is weak \emph{because it is global}---it feels the entire CPT-symmetric Universe, not just our sector. This structural dilution replaces the usual hierarchy problem with a simple projection effect.

\subsection{The Graviton as Global Propagator}
The graviton is the quantum of the single shared metric $g_{\mu\nu}$. In RDCC it is therefore \emph{not} a sectoral excitation but the carrier of global correlations between $\mathcal{H}_+$ and $\mathcal{H}_-$. At the CPT fixpoint (the cosmological bounce) the conditions
\[
H(t_b) = 0, \quad \alpha(t_b) \to 0, \quad \rho_{\rm tot} = \rho_{\rm crit}
\]
make the two sectors maximally entangled. The graviton at this point propagates the \emph{pure global CPT eigenstate}. There is no need for additional UV degrees of freedom or wormholes: the graviton already \emph{is} the entanglement structure of the two sectors.

\subsection{Summary}

RDCC does not attempt to replace quantum mechanics or quantum gravity. Instead, it provides a structural framework in which many microphysical puzzles appear as artifacts of sectoral projection. Time, entropy, nonlocality, collapse, causal limits, and information loss are not fundamental inconsistencies but perspectival effects arising from the CPT-symmetric global state. This perspective does not solve these problems dynamically but clarifies their conceptual origin.

\section{Conclusion}

The structural principles of RDCC provide a unified interpretative framework in which many conceptual features of cosmology, quantum theory, and gravitation appear as perspectival effects of sectoral projection rather than fundamental inconsistencies. Global CPT symmetry, non-separability, and the existence of a structural fixpoint together define the boundary conditions from which sectoral notions of time, entropy, causality, and locality emerge.

The interpretative extensions developed in this companion document illustrate how phenomena such as black hole interiors, the speed of light, high-curvature regimes, holography, timeless quantum cosmology, decoherence, teleportation, ER=EPR correlations, Many-Worlds and black-hole interiors can be understood without introducing new fields or modifying the core RDCC dynamics. These ideas clarify the structural coherence of the two-sector framework and highlight the limits of sectoral description.

The final section has shown that several long-standing microphysical puzzles—including the problem of time, the thermodynamic arrow, Bell nonlocality, the measurement problem, the weakness of gravity, and the nature of dark matter—can be reframed within the RDCC perspective. While no dynamical solutions are proposed, the structural reinterpretation offered here provides conceptual clarity and a coherent informational viewpoint.

This document is therefore intended as a conceptual companion to the RDCC framework: it expands the interpretative landscape, deepens the structural understanding of CPT-symmetric cycles, and preserves the minimalist and non-competitive philosophy that defines RDCC.

\bibliographystyle{apsrev4-2}
\nocite{*}
\bibliography{references}

\end{document}